\scanpage{24}
\mspage{20}
% archive sheet number: 23

``\unsure{magasin} résiduel'' ou ``complémentaire'' p.r. à une des extrémités de $I$, ou des deux. Supposons p.ex. \deleted{qu'il n'y ait} qu'on regarde le $(I,\partial I)\setminus\{\omega_0\}$ (façon de parler) -- ou, plus précisément, le sous-magasin
\[
\deleted{\unsure{M_{\overset{\circ}{\ll}}}}\operatorname{Omb}^{\circ}(\{I_{\omega_0,\omega_1},\omega_1\})=\unsure{M|\mathcal M}\setminus\omega_0
\]
(partie de $M$ fermée pour $\unsure{\omega_0|\overset{\circ}{\ll}}$), on désire l'interpréter en termes de
\[
(\overbrace{I\setminus\{\omega_0\}}^J,\omega_1)
\]
un ordonné $J$, \uline{muni} d'un plus gd élément $\omega_1$. On garde toutes les strates de $M(I,\partial I)$, à l'exception de $\omega_0$. Mais les strates de la forme \inserted{avec $\omega_0<a$ i.e. $a\in J$} $I_{\omega_0,a}$, \deleted{s'écrivent} maintenant
\[
J_{<a}
\]
\margin{On ne suppose plus ici que $a$ ne soit pas \unsure{minimal}, cf plus bas}
Les relations $\unsure{\trianglelefteq|\triangleleft},\leq,\overset{\circ}{\ll},\mathrel{|\circ|}$ \deleted{et} dans l'ens.
\[
M(J,\omega_1)\deleted{\unsure{=}}\quad(=M(I,\partial I\setminus\{\omega_0\})
\]
pourront alors s'expliciter directement, en termes de $J,\omega_1$, sans réf. \inserted{explicite} à $(I,\partial I)$, en travaillant avec les strates
\[
a,\ J_{<a},\ \unsure{\mathbb I|I|J}_{a,b}\qquad a<b,\quad a,b\in J.
\]
Si p.ex. $I=\{\omega_0,\omega_1\}$, on aura $J=\{\omega_1\}$,
